\documentclass[12pt]{article}
\usepackage[utf8]{inputenc}
\usepackage[T1]{fontenc}
\usepackage{hyperref}
\title{T’en Language Resources for K-12 Curriculum Development}
\author{Compiled by Claude Sonnet 5 (claude-sonnet-5) for Andr\'as Kornai}
\date{}
\begin{document}
\maketitle
\section{Introduction}
T’en (Then, Rao, Rau, Yanghuang) is a Kam-Sui language of Guizhou, China (east Pingtang, Dushan and Huishui counties, per Verbix); about 20,000 speakers (ASJP). Glottolog places it in Tai-Kadai, Kam-Sui; some scholars regard it as a Sui dialect. Glottocode tenn1245. The Glottolog assessment is that it is shifting and threatened.
\section{Primary Educational Resources}
None found.
\section{Linguistic and Reference Materials}
\begin{itemize}
\item Glottolog entry: \url{https://glottolog.org/resource/languoid/id/tenn1245}
\item ASJP word list: \url{https://asjp.clld.org/languages/TEN}
\item Verbix language page: \url{https://docs.verbix.com/Languages/Ten}
\end{itemize}
\section{Religious and Community Resources}
\begin{itemize}
\item Joshua Project language page: \url{https://joshuaproject.net/index.php/languages/tct}
\end{itemize}
\section{Digital and Audio Resources}
None found.
\section{Curriculum Development Context}
Only catalogue entries and a word list surfaced. No teaching materials, orthography or school programmes were found. Materials for the closely related Sui language may be a possible starting point, but were not searched for here.
\end{document}